
%! Author = lili
%! Date = 4/30/26

de bruijn subgraph
DBG(s,k)  = (V,E)

V 0 alle k-mere von s

Wie viele k-mere hat eine Sequenz der Länge n sollte man wissen: n-k+1

info: für jeden DB Graph gibt es immer einen Eulerpfad

\begin{theo}
    Theorem 2.1 auf deutsch: \# worte mit gleichen (k+1)-mer Profil wie s = \# Euler-Pfade in DBG(s,k)
\end{theo}

\begin{kom}
    man könnte bei bsp 2.1. denken hey bei a kann es nur einen EP geben, weil es kein Kreis ist, und bei 2b
    mus es dann
    mehrere geben, eil da eine Schleife drin ist, abre NEIN der Pfad ist die Reihenfolge der Knoten und nicht der
    Kanten, also gibt es nur einen (visualisieren) Man kann nicht unterscheiden, ob man zuerst die obere oder untere
    Kante der Schleife zuerst gegangen ist. Gibt also nur eine sequenz mit em k+1 meer profil. Denn die Seuenz wird
    ja aus den Knoten gemacht, nicht den Kanten...sozusagen.
\end{kom}